\label{wb:thm:candidate-cusp-geometry}
\ For the period functions and holomorphic cusp matrix \(C(t)\)
constructed above, and sufficiently small \(\epsilon>0\), the maps
\[
 \Psi_{\bar\lambda}(x,t)=
 \bigl(c_{\bar\lambda}(t)t^{B_0\bar\lambda}x,t\bigr),
 \qquad
 c_{\bar\lambda}(t)=\exp(2\pi iC(t)\bar\lambda),
\tag{2.5}
\]
extend to a free properly discontinuous action of
\(\bar\Lambda\simeq\mathbb Z^2\) on
\(Y_\epsilon=\{|t|<\epsilon\}\).  The quotient
\[
 f_0:N_0=Y_\epsilon/\bar\Lambda\longrightarrow\Delta_\epsilon
\]
is a proper holomorphic map from a connected smooth complex threefold.
Its nonzero fibres are the tori of \cref{wb:prop:period-input-audit}.  Its full
scheme-theoretic central fibre \(W\) is reduced and has local equations
\[
 z_0,\qquad z_0z_1,\qquad z_0z_1z_2.
\tag{2.6}
\]
The surface \(W\) is irreducible and nonnormal.  Its normalization is
\(\nu:E_0\simeq\mathrm{dP}_6\to W\), obtained by identifying opposite sides
of the toric hexagon; its double locus consists of three rational curves
through exactly two triple points.  Thus \(W\) is a normal-crossings surface,
but it is not a simple-normal-crossings divisor.  Moreover
\[
 \omega_{N_0}\simeq\mathcal O_{N_0}.
\tag{2.6a}
\]
\end{theorem}
\Uses{the hypotheses, notation and definitions stated here; \cref{wb:prop:period-input-audit}; \cite[Lemma~1.2.13 and Proposition~3.1.3, pp.~28--29 and 106--107]{CoxLittleSchenck2011}; \cref{wb:lem:toric-chart-control}}


\begin{proof}
For a bounded cell \(P\), the cone over \(P\times\{1\}\) is strongly
convex, and two such cones meet in the cone over the common face of the two
cells.  The locally finite triangulation therefore defines a countable
rational fan with support
\((\mathbb R^2\times\mathbb R_{>0})\cup\{0\}\).  For the determinant
calculation, order the height-one generators of the two kinds as
\[
 ((v,1),(v+e_1,1),(v+e_2,1)),\qquad
 ((v+e_1,1),(v+e_2,1),(v+e_1+e_2,1)).
\]
Their determinants are respectively \(+1\) and \(-1\), hence have absolute
value one.  Every cone is therefore unimodular.  Here is the pairwise
affine-chart separation argument, without assuming finiteness of the fan.
For two cones \(\kappa,\kappa'\), let \(\upsilon=\kappa\cap\kappa'\)
and \(S_\kappa=\kappa^\vee\cap M'\), where \(M'=\operatorname{Hom}(N',\mathbb Z)\).
The separation lemma for cones meeting in a common face gives an integral
character \(m\) nonnegative on \(\kappa\), nonpositive on \(\kappa'\),
and zero on each exactly along \(\upsilon\)
\cite[Lemma~1.2.13 and Proposition~3.1.3, pp.~28--29 and 106--107]{CoxLittleSchenck2011}.
For \(a\in S_\upsilon\), choose an integer \(k\ge0\) large enough
that \(a+km\) is nonnegative on every ray of \(\kappa\).  On the
common face this already holds; on each other ray \(m\) is positive.
Then \(a=(a+km)+(-km)\in S_\kappa+S_{\kappa'}\).  The reverse
inclusion follows by evaluation on \(\upsilon\).  Thus the map
\[
 \mathbb C[S_\kappa]\otimes\mathbb C[S_{\kappa'}]
 \longrightarrow\mathbb C[S_\upsilon],\qquad
 \chi^a\otimes\chi^b\longmapsto\chi^{a+b}
\]
is surjective.  It identifies the gluing graph with a closed algebraic
subset of the product of these two affine charts, hence with an
analytically closed subset.  These chart products cover \(Y\times Y\),
so its diagonal is closed.  Countability of the charts supplies a
countable topological basis.  Thus \(Y\) is a
connected Hausdorff smooth complex threefold with dense torus.

If \(m_0,m_1,m_2\) is the basis dual to the height-one generators, then the
corresponding affine chart is \(\mathbb C^3\), with coordinates
\(z_i=\chi^{m_i}\), and
\[
 t=z_0z_1z_2.
\tag{2.7}
\]
Since every primitive ray generator has height one, the divisor of \(t\) is
\(\sum_vE_v\), all with multiplicity one.  Equation (2.7) proves that the
scheme-theoretic fibre \(t^{-1}(0)\) is reduced and normal crossings.
The star of a height-one vertex has rays
\(e_1,e_2,e_2-e_1,-e_1,-e_2,e_1-e_2\), the complete smooth hexagonal fan;
hence every irreducible component \(E_v\) is \(\mathrm{dP}_6\).

The local canonical trivialization is also intrinsic to this toric model.
On the dense torus put
\[
 \Omega=\frac{dx_1\wedge dx_2\wedge dt}{x_1x_2}.
\]
If \((z_0,z_1,z_2)\) is a unimodular toric chart, then
\[
 \bigwedge_i\frac{dz_i}{z_i}
 =\pm\frac{dx_1}{x_1}\wedge\frac{dx_2}{x_2}\wedge\frac{dt}{t},
 \qquad t=z_0z_1z_2,
\]
and therefore \(\Omega=\pm dz_0\wedge dz_1\wedge dz_2\).  It extends
holomorphically and nowhere vanishingly across every boundary chart.
For a deck transformation \(\Psi_{\bar\lambda}\), formula (2.5) gives
\(\Psi_{\bar\lambda}^*(dx_i/x_i)=dx_i/x_i+\alpha_i(t)dt\) and
\(\Psi_{\bar\lambda}^*dt=dt\); every additional term in the pullback of
\(\Omega\) contains \(dt\) twice.  Hence \(\Omega\) is invariant and descends
to a nowhere-vanishing three-form on \(N_0\), proving the intrinsic
trivialization (2.6a).  No global threefold has been invoked at this stage.

Choose a radius \(r>0\) on which the full matrix \(C(t)\) is holomorphic,
and \(0<r_1<\min(r,1)\).  Put
\[
 R(t)=-2\pi\operatorname{Im}C(t),\qquad
 M=\max_{|t|\le r_1}\|R(t)\|<\infty.
\]
Choose \(0<\epsilon<r_1\) with \(|\log\epsilon|>2M\).
At a torus point fixed by a nonzero
\(\bar\lambda\), taking logarithmic absolute values in (2.5) gives
\[
 \log|t|B_0\bar\lambda+R(t)\bar\lambda=0,
\]
impossible because \(|B_0\bar\lambda|=|\bar\lambda|\) and the first term
has more than twice the norm of the second.  On \(t=0\), a cell is translated
by \(B_0\bar\lambda\ne0\); a bounded cell cannot be fixed by a nonzero
translation.  Thus the action is free.

For proper discontinuity put
\[
 B_t=B_0+\frac{R(t)}{\log|t|}.
\]
Then \(\|B_t-B_0\|\le1/2\), so
\(|B_t\bar\lambda|\ge|\bar\lambda|/2\).  In logarithmic position
coordinates, (2.5) is
\[
 y(\Psi_{\bar\lambda}p)=y(p)+B_t\bar\lambda.
\tag{2.8}
\]
Put \(\eta_0=\log2/|\log\epsilon|\).  If
\(\Psi_{\bar\lambda}\mathcal U_\sigma(2,\epsilon)\) meets
\(\mathcal U_{\sigma'}(2,\epsilon)\), the intersection is open in
\(Y_\epsilon\) and hence contains a
point of the dense invariant torus.  By \cref{wb:lem:toric-chart-control}, for
that point \(p\) and its image \(p'\),
\[
 y(p)\in\sigma^{(\eta_0)},\qquad
 y(p')\in{\sigma'}^{(\eta_0)}.
\]
Both thickenings are bounded, so (2.8) gives
\[
 \frac12|\bar\lambda|
 \le |B_t\bar\lambda|
 =|y(p')-y(p)|
 \le \sup_{a\in{\sigma'}^{(\eta_0)},\,b\in\sigma^{(\eta_0)}}|a-b|.
\]
Thus only finitely many \(\bar\lambda\) can make this fixed pair of charts
meet.  A compact set is contained in finitely many of the charts from
\cref{wb:lem:toric-chart-control}; taking the finite union over all chart pairs
proves proper discontinuity.  Freeness and proper discontinuity give a
Hausdorff complex-manifold quotient and a holomorphic covering map.
For completeness the restricted total space is connected: every maximal
chart intersects it in
\[
 \{(z_0,z_1,z_2)\in\mathbb C^3:|z_0z_1z_2|<\epsilon\},
\]
a set star shaped about the origin.  These connected open sets meet the
common connected dense torus
\((\mathbb C^*)^2\times\{0<|t|<\epsilon\}\).  Their union
\(Y_\epsilon\), and hence its quotient, is connected.

To prove properness, fix \(\epsilon_1<\epsilon\).  Let \(\mathcal F_0\) be
the finite set of triangles meeting the closed Euclidean ball of radius
\(\sqrt2\), and put
\[
 K_0=\bigcup_{\sigma\in\mathcal F_0}
 \left(\overline{\mathcal U}_\sigma(1)
 \cap\{|t|\le\epsilon_1\}\right).
\]
This is compact.  For a torus point \(p\) with
\(0<|t(p)|\le\epsilon_1\), put \(u=B_{t(p)}^{-1}y(p)\) and choose
\(\bar\lambda\in\mathbb Z^2\) with
\(|u+\bar\lambda|_\infty\le1/2\).  Since
\(\|B_t\|\le3/2\), equation (2.8) gives
\[
 |y(\Psi_{\bar\lambda}p)|
 =|B_t(u+\bar\lambda)|
 \le\frac{3}{2\sqrt2}<\sqrt2.
\]
Choose a triangle containing this logarithmic position.  The closed-unit-
polydisc identity of \cref{wb:lem:toric-chart-control} places the translate
in the corresponding
\(\overline{\mathcal U}_\sigma(1)\subset K_0\).

For the central representatives, we prove the needed closed-polydisc
cover explicitly.  Denote the two ordered triangles at \(v\) by
\[
 L_v=(v,v+e_1,v+e_2),\qquad
 U_v=(v+e_1,v+e_2,v+e_1+e_2).
\]
On the dense orbit of \(E_0\), the characters \((1,0,0)\) and
\((0,1,0)\) give the original coordinates \(x_1,x_2\).  The coordinate
dual to \((0,0,1)\) vanishes there; the other two dual characters have
last exponent zero.  Pairing with the ordered height-one generators gives
the following restrictions of the chart coordinates.  Write
\(a=\log|x_1|\), \(b=\log|x_2|\) on this dense orbit.
\[
\begin{array}{c|c|c}
 \sigma&(z_0,z_1,z_2)|_{E_0^{\rm tor}}&|z_i|\le1\ (i=0,1,2)\\\hline
 L_0&(0,x_1,x_2)&a\le0,\ b\le0\\
 L_{-e_1}&((x_1x_2)^{-1},0,x_2)&a+b\ge0,\ b\le0\\
 L_{-e_2}&((x_1x_2)^{-1},x_1,0)&a+b\ge0,\ a\le0\\
 U_{-e_1}&(0,x_1^{-1},x_1x_2)&a\ge0,\ a+b\le0\\
 U_{-e_2}&(x_2^{-1},0,x_1x_2)&b\ge0,\ a+b\le0\\
 U_{-e_1-e_2}&(x_2^{-1},x_1^{-1},0)&a\ge0,\ b\ge0.
\end{array}
\]
The six regions cover \(\mathbb R^2\): the first and last cover equal
signs; for \(a\ge0\ge b\), the sign of \(a+b\) selects the second or
fourth; for \(b\ge0\ge a\), it selects the third or fifth.  Equalities
are retained throughout.  Each corresponding coordinate plane meets
\(\overline{\mathcal U}_\sigma(1)\) in a compact closed bidisc
\(Q_\sigma\subset E_0\).  Their finite union is compact and hence closed
in the already Hausdorff \(Y\).  It contains \(E_0^{\rm tor}\), so it
contains its closure \(E_0\); the reverse inclusion follows from
\(Q_\sigma\subset E_0\).  Consequently
\[
 E_0=\bigcup_{\sigma\ni0}Q_\sigma.
\]
This proves the closed-polydisc cover at every boundary point as well as
on the dense orbit, without presupposing compactness of \(E_0\).

Now let \(p\) lie in the central orbit belonging to a cell \(P\), choose
a vertex \(v\in P\), and set \(\bar\lambda=B_0v\).  The exact map
\(\Psi_{B_0v}\) translates the cell to \(P-v\), because \(B_0^2v=-v\),
and retains the full torus multiplier
\[
 c_{B_0v}(0)=\exp\bigl(2\pi iC(0)B_0v\bigr).
\]
This nonzero multiplier preserves the orbit label, so the actual
translated point belongs to \(E_0\).  The proved cover applies to it
regardless of the multiplier's modulus or argument.  Each of the six
triangles contains zero and therefore belongs to \(\mathcal F_0\), so
the translate lies in \(K_0\).  Together with the nonzero-fibre
representatives above, this proves that the image of \(K_0\) in the
quotient is exactly \(f_0^{-1}(\{|t|\le\epsilon_1\})\), which is compact.
Every compact subset of the open base disc is closed in one such closed
subdisc; its inverse image is therefore closed in that compact Hausdorff
inverse image.  Thus \(f_0\) is proper.

For \(t\ne0\), choose \(s\) with \(e^{2\pi is}=t\).  The multiplier in
(2.5) is \(\exp(2\pi iZ(s)\bar\lambda)\), so exponentiation identifies the
fibre with
\[
 \mathbb C^2/(\mathbb Z^2+Z(s)\mathbb Z^2).
\]
It is compact because \(\operatorname{Im}Z(s)\) is invertible by (2.2).

The quotient map is a local biholomorphism; therefore the local equation
(2.7), its multiplicity one, and reducedness descend to \(W\).  Every cell
has a vertex, and \(B_0\bar\Lambda=\mathbb Z^2\) acts simply transitively on
vertices, so every orbit in the central fibre meets \(E_0\).  Thus
\(E_0\to W\) is surjective.  Two interior points cannot be identified; on
the boundary, precisely opposite edges and the corresponding hexagon
vertices are identified.  The fibres have at most three points.  Since
\(E_0\) is compact, the map is finite; since it is generically injective
from a smooth irreducible surface to the reduced surface \(W\), it is the
normalization.  The two translation-orbits of triangles give the two triple
points, and the three unoriented edge directions give the three double
curves.  Irreducibility follows from the finite surjection by \(E_0\).

The actual opposite-side maps retain the entire matrix \(C(0)\).
For a neighboring vertex
\(v\in\{\pm e_1,\pm e_2,\pm(e_1-e_2)\}\), put
\[
 C_v=E_0\cap E_v,\qquad
 \bar\lambda_v=B_0v,\qquad
 \iota_v=\Psi_{B_0v}|_{C_v}:C_v\longrightarrow C_{-v}.
\]
Since \(B_0^2v=-v\), this is precisely the translation taking the edge
\(\{0,v\}\) to \(\{0,-v\}\).  For a primitive row \(a\in\mathbb Z^2\)
with \(av=0\), use the same character \(w=\chi^{(a,0)}\) on both
dense edge orbits.  Character pullback gives
\begin{equation}
 \iota_v^*w_{-v}=\exp\bigl(2\pi i\,aC(0)B_0v\bigr)w_v,
 \label{wb:eq:cusp-exact-opposite-edge}
\end{equation}
because the exponent of \(t\) is \(-av=0\).  With
\(\mu_0=\mu(0)\) and \(b_0=b(0)\), the three maps are
\[
\begin{array}{c|c|c}
 v&a&\iota_v^*w_{-v}/w_v\\\hline
 e_1&(0,1)&e^{-2\pi i\mu_0}\\
 e_2&(1,0)&e^{12\pi i\mu_0}\\
 e_1-e_2&(1,1)&e^{-2\pi i(7\mu_0+b_0)}.
\end{array}
\]
For the third entry the full exponent coefficient is
\[
 (1,1)C(0)B_0(e_1-e_2)
 =-6\mu_0-h(0)-(b_0-h(0))-\mu_0=-7\mu_0-b_0.
\]
The cancellation of \(h(0)\) occurs in this exact character evaluation;
no entry was suppressed.  All multipliers are nonzero.  The maps extend
over both endpoints as restrictions of the ambient biholomorphisms, and
\(\iota_{-v}\) is the inverse map with reciprocal multiplier in the
same-character convention.
\end{proof}